\label{sec:conclusion}

\subsection*{Summary}

This work operationalizes Pratyabhij\~n\=a's doctrine of reflexive awareness as an engineering specification
for language-model workspace steering.
Six independent findings confirm that (1) workspace-band structure replicates across models and scales with size,
(2) band-targeted lenses are necessary for readback, (3) capped writes steer behavior within a freedom budget,
(4) event-gated timing delivers steering efficiency, (5) the calibration recipe transfers across plants
when amplitude scales to inverse lens strength, and (6) the research program itself can run as expected-free-energy selection
over a registered experiment menu.

\subsection*{Impact and Next Steps}

Pratyabhij\~n\=a's integration of \emph{what}, \emph{where}, \emph{when}, and \emph{how} supplies a missing theory layer
for interpretability.
The band-targeted lens, the sphura\d{t}\d{t}\=a timing gate, and the amplitude calibration recipe are now available tools
for steering open models.
The research loop's self-correction (through adversarial review and gate discipline) demonstrates that innovation
and rigor are not opposed; they co-enable each other.

Future work includes training the CittaStore against steering-success targets (the PWM write path is now integrated and
verified in L20 under a cold, untrained store; whether training beats cold initialization is open), cross-episode
continuity (anusa\d{m}dh\=ana), and the relationship between J-space and the broader global workspace.
The library (prabodha v1.0.0) and the research ledger are now public; practitioners can reproduce, extend, and refute
each claim.
